\paragraph{Step 1 (T0) structure.}
\begin{itemize}[nosep,leftmargin=*]
\item The first MODEL of a multi-model file defines the atoms. Later models serve as trajectory frames.
\item For alternate locations, the first conformer (A, 1 or blank) is kept.
\item Virtual sites are removed on reading. Examples are the $\sigma$-hole extra point (EP) of a halogen, LP sites and TIP4P M sites. DCD and XYZ frames with these sites are mapped onto the remaining atoms.
\item Bonds come from distances: d $<$ r\_cov(1) + r\_cov(2) + 0.45 \AA{}. Links between residues are limited to peptide, disulfide and phosphodiester bonds.
\item Element symbols come from columns 77--78. Atom names supply missing element symbols.
\item Residues are classed as protein, nucleic acid, water, ion (single atom) or ligand (all other residues).
\end{itemize}
\paragraph{Step 2 (T1) geometry.}
\begin{itemize}[nosep,leftmargin=*]
\item Donors: C-bound Cl, Br and I atoms of ligands and modified residues.
\item A halogen bond needs d(X$\cdots$Y) $\le$ the pair cut-off and 140$^\circ$ $\le$ $\theta_1$ $\le$ 180$^\circ$.
\item Acceptors: protein N/O/S, nucleic-acid N/O, aromatic rings (Phe, Tyr, Trp) with $\alpha$ $\le$ 60$^\circ$.
\item R = d / (r\_X + r\_Y) with Bondi radii. An aromatic centroid has r = 1.77 \AA{}.
\item G = [g\_d $\cdot$ g\_$\theta$ $\cdot$ ($\frac12$ + $\frac12$ g$_2$)]\^{}$\frac13$. g\_d = 1 at R $\le$ 0.85. g$_2$ is a Gaussian around $\theta_2$ = 120$^\circ$ (width 40$^\circ$).
\item Poor Lewis bases are kept with the flag ``weak''. The NBO gate decides their fate.
\item Hydrogens are not required. Periodic images apply only to MD boxes (orthorhombic, edges $\ge$ 10 \AA{}).
\end{itemize}
\paragraph{Step 3 (T2) dynamics.}
\begin{itemize}[nosep,leftmargin=*]
\item Frames must have the atoms of the topology in the same order.
\item The T1 test runs on every frame. Contacts in fewer than 10\% of the frames are not reported.
\item Residence time $\tau$ = mean length of unbroken runs. D = f$\cdot$(1 $-$ e\^{}($-\tau$/5))\^{}0.25.
\item The representative frame has the shortest d. Ties go to the largest $\theta_1$.
\item Orthorhombic boxes (DCD unit cell or CRYST1) use the minimum-image convention. The ligand should be imaged and centred.
\end{itemize}
\paragraph{Step 4 (T3a) QM cluster model.}
\begin{itemize}[nosep,leftmargin=*]
\item Protonation states and partial charges come from AMBER ff14SB (+ OL15 DNA, OL3 RNA). ASP and GLU are anionic. LYS and ARG are cationic. CYS and TYR are neutral. HIS is HIE by default. A tautomer name in the file (HID/HIE/HIP, HSD/HSE/HSP) or its hydrogens override this default. Changes for single residues override everything.
\item Cys pairs with S--S $<$ 2.5 \AA{} are disulfides (CYX). Chain termini are charged (NH$_3^+$, COO$^-$). Residues next to a chain break are treated as internal.
\item Fragment 1 of the QM region is the whole halogenated ligand. Fragment 2 holds the whole acceptor residues.
\item Missing hydrogens follow the force-field template of each residue (number of H per atom) with VSEPR geometry. Cut bonds get H caps (fixed X--H length).
\item Ligands: one heavy-atom structure fixes the hydrogens and the charge. Every sp$^2$ carbon gets one double bond. The keto form is preferred. An sp$^2$ ring N without a double bond carries one H. P keeps one P=O bond. S keeps up to two S=O bonds. At pH 7, COOH, phosphate, phosphonate, sulfonate and sulfate groups lose their O--H hydrogen. Amines stay neutral. Hydrogens of the file stay only with a matching number. The PDB formula (FORMUL) checks the count.
\item Fragment charges: residue formal charges come from the force-field template. The ligand charge comes from a loaded MOL2 file or from the ligand structure above. Every value can be overridden.
\item The QM cluster sits in an implicit dielectric $\varepsilon$ = 4 (CPCM/PCM). The dielectric represents the protein.
\end{itemize}
\paragraph{Step 4 (T3a) QM/MM model (extra lines).}
\begin{itemize}[nosep,leftmargin=*]
\item QM/MM: MM point charges within 15 \AA{} (whole residues) come from a loaded PQR file or from AMBER ff14SB (+ OL15 DNA, OL3 RNA). Water model: TIP3P.
\item For atoms without hydrogens in the file, each hydrogen charge moves onto its heavy atom. Residue charges stay exact. Local dipoles are approximate.
\item Boundary: the charge of each M1 atom moves to the midpoints of its M1--M2 bonds (redistributed charge, RC).
\item The embedding is electrostatic only. QM--MM van der Waals terms are not included. They do not enter the NBO analysis of a single point.
\item Ligands and cofactors in the MM region carry MOL2 charges from a loaded file. Without a MOL2 file, a charged group (carboxylate, phosphate, sulfate, R$_4$N$^+$) carries its formal charge on its heavy atoms. Neutral groups get zero charges. Both cases are flagged.
\end{itemize}
\paragraph{Step 5 (T3b) NBO gate.}
\begin{itemize}[nosep,leftmargin=*]
\item A contact is verified at $\Sigma$E(2) $\ge$ 0.5 kcal/mol. The sum covers n(Y)$\to\sigma$*(A--X) and n(Y)$\to$RY*(X) for lone-pair acceptors. It covers $\pi$(ring)$\to\sigma$*(A--X) for aromatic acceptors.
\item The output geometry must match the manifest atom by atom (same elements, deviation $<$ 0.05 \AA{}). A mismatch blocks the gate.
\item NBO prints only couplings above its own print threshold (0.50 kcal/mol by default). Open-shell $\alpha$ and $\beta$ tables are summed.
\end{itemize}
\paragraph{Step 6 (T4) ALMO-EDA.}
\begin{itemize}[nosep,leftmargin=*]
\item Fragment 1 is the halogenated ligand. Fragment 2 is the acceptor environment. Both keep the charges and multiplicities of the QM model.
\item $\chi$\_CT = E\_CT / (E\_POL + E\_CT). Values $\ge$ 0.6 mean charge-transfer driven. Values $\le$ 0.4 mean polarization driven. Values in between mean mixed.
\item Energies are given in kJ/mol (1 kcal = 4.184 kJ).
\end{itemize}
\paragraph{Step 7 (T5) BXQI.}
\begin{itemize}[nosep,leftmargin=*]
\item BXQI = 100 $\cdot$ $\Pi$ s\_k\^{}(w\_k/$\Sigma$w) over all completed tiers. Weights: G 0.25, D 0.2, E 0.3, C 0.25.
\item s\_E = 1 $-$ e\^{}($-\Sigma$E(2)/2). s\_C = 1 $-$ e\^{}($-$|$\Delta$E\_int|/10). Grades: A $\ge$ 75, B $\ge$ 50, C $\ge$ 25, else D.
\item Grade A requires a passed NBO gate. A contact with a failed gate receives grade C or D.
\item An NBO/ALMO inconsistency is flagged in two cases: E(2) $>$ 4 kJ/mol with $-$E\_CT $<$ 0.5 kJ/mol, and E(2) $<$ 1 kJ/mol with $-$E\_CT $>$ 4 kJ/mol.
\end{itemize}
